\clearpage
\begingroup
\HistoricalHead{237}{R}{WEITERES ZUR BAHNBESTIMMUNG.}
\begin{center}
[3.]\par\medskip
\textls[100]{Vorschriften,}\\
um aus\\
Eccentricität $=e=\sin\varphi$, und wahrer Anomalie $=M$\\
\textls[100]{abzuleiten}\\
eccentrische Anomalie $=E$, und mittlere Anomalie $=N$
\end{center}
\[
\begin{array}{rl}
\mathrm{I.}&\operatorname{tang}\tfrac12E=\operatorname{tang}\tfrac12M\operatorname{tang}(45^0+\tfrac12\varphi),\\
\mathrm{II.}&N=E+\sin E\sin\varphi.
\end{array}
\]

Ist zugleich der Abstand $r$ bekannt und $a,b,p$, so bedient man sich öfters mit Vortheil folgender Formeln:
\[
\sin E=\sin M\mathbin{.}\frac rb,
\]
\[
\cos E=\frac ra\cos M-e=\frac rp(\cos M-e)=\frac{\cos M-e}{1-e\cos M}.
\]
\noindent{\small [S. 12]}\par\noindent Auch folgende Formeln sind sehr brauchbar:
\[
\begin{array}{rl}
\mathrm{I.}&\operatorname{cotg}\tfrac12(E-M)=\dfrac{\operatorname{cotg}\tfrac12\varphi}{\sin M}-\operatorname{cotg}M,
\end{array}
\]
\clearpage
\HistoricalHead{238}{L}{THEORETISCHE ASTRONOMIE. NACHLASS.}
\SourceLeaf{243}
\[
\begin{array}{rl}
\mathrm{II.}&\operatorname{cotg}\tfrac12(E'-E)=\dfrac{\operatorname{cotg}\tfrac12(M'-M)}{\cos\varphi}
-\dfrac{\operatorname{tang}\varphi\cos\tfrac12(M'+M)}{\sin\tfrac12(M'-M)},\\[6pt]
\mathrm{III.}&\sin(E'-E)=\dfrac{rr'}{ab}\sin(M'-M)\mathbin{.}\{-N'+N+E'-E\}.
\end{array}
\]
\begin{center}\textls[100]{Beispiele.}\end{center}
\noindent I.
\[
\begin{array}{rl}
v=&69^0\,14'\,22''77\\
\omega=&48\;30\;14''05\\\cline{1-2}
M=&20^0\,44'\,8''72\\[4pt]
\tfrac12M=&10\;22\;4''36\\
\log\operatorname{tg}\tfrac12M=&9{,}262\,3440\\
\log\operatorname{tg}(45^0+\tfrac12\varphi)=&0{,}080\,5506\\\cline{1-2}
{}[\log\operatorname{tg}\tfrac12E]=&9{,}342\,8946\\
\tfrac12E=&12^0\,25'\,13''763\\
E=&24\;50\;27''526\\
&+\;4\;24\;59''978\\\cline{2-2}
N=&29^0\,15'\,27''50
\end{array}
\]
\[
\begin{array}{rl}
\varphi=&10^0\,33'\,59''43\\
45^0+\tfrac12\varphi=&50\;16\;59''715\\[6pt]
\log\sin E=&9{,}623\,3539\\
\log\sin\varphi=&9{,}263\,3443\\\cline{1-2}
&8{,}886\,6982\\
\log1''=&4{,}685\,5749\\\cline{1-2}
&4{,}201\,1233
\end{array}
\]
\[
15\,899''978=4^0\,24'\,59''978\text{\textsuperscript{[*)]}}
\]
\begingroup\renewcommand{\thefootnote}{[*)}\footnotetext{Dieser Wert beruht auf einem Versehen; er sollte lauten: $15\,889''978=4^0\,24'\,49''978$. Er wurde nicht verbessert, weil damit das Beispiel in der folgenden Nr. [4.] gerechnet worden ist.]}\endgroup
\begin{center}Contrôle\end{center}
\[
\begin{array}{rl}
\log\sin M=&9{,}549\,0749\\
\log r=&0{,}425\,7296\\
\mathrm{C.}\log b=&9{,}648\,5494\\\cline{1-2}
{}[\log\sin E=]&9{,}623\,3539
\end{array}
\]
\noindent II.
\[
\begin{array}{rl@{\qquad}l@{\qquad}rl}
v=&73^0\,48'\,36''17&&E=&30^0\,14'\,45''012\\
\omega=&48\;30\;14''05&\text{[die Rechnung ergab]}&N=&35^0\,32'\,17''47\\\cline{1-2}
M=&25^0\,18'\,22''12&&&
\end{array}
\]
\noindent{\small [S. 13]}\par
\noindent III.
\[
\begin{array}{rl@{\qquad}l@{\qquad}rl}
v=&78^0\,26'\,32''49&&E=&35^0\,40'\,56''998\\
\omega=&48\;30\;14''05&\text{[die Rechnung ergab]}&N=&41^0\,48'\,39''71\\\cline{1-2}
M=&29^0\,56'\,18''44.&&&
\end{array}
\]
\begin{center}
[4.]\par\medskip
\textls[100]{Vorschrift}\\
zur Bestimmung der Zeit, in welcher eine gegebne mittlere Bewegung geschieht.
\end{center}
\noindent\hspace*{\parindent}Diese erhält man in Tagen, wenn man die mittlere Bewegung mit $a^{3/2}$ multiplicirt, und das Product mit der tägl[ichen] mittlern Bewegung der Sonne dividirt. Diese ist $3548''19223$, $\log=3{,}550\,0071$.
\clearpage
\HistoricalHead{239}{R}{WEITERES ZUR BAHNBESTIMMUNG.}
\SourceLeaf{244}
\begin{center}\textls[100]{Beispiele.}\end{center}
\[
\begin{array}{rl}
N=&29^0\,15'\,27''50\\
N'=&35\;32\;17''47\\
N''=&41\;48\;39''71
\end{array}
\qquad
\begin{array}{rl}
N'-N=&6^0\,16'\,49''97=22\,609''97\\
N''-N'=&6\;16\;22''24=22\,582''24
\end{array}
\]
\[
\begin{array}{rl}
\log a=&0{,}358\,8782\\
{}[\tfrac32\log a=&0{,}538\,3173]\\
{}[\mathrm{C.}\log3548''\;\text{. . .}=]&6{,}449\,9929\\\cline{1-2}
&6{,}988\,3102
\end{array}
\]
\[
\begin{array}{rl@{\qquad}rl}
\log(N'-N)=&4{,}354\,3000&\log(N''-N')=&4{,}353\,7670\\
&6{,}988\,3102&&6{,}988\,3102\\\cline{2-2}\cline{4-4}
&1{,}342\,6102&&1{,}342\,0772\\[4pt]
&22^{\mathrm d}\,009\,5025&&21^{\mathrm d}\,982\,505
\end{array}
\]
\begin{center}[5.]\end{center}
\noindent{\small [S. 14]}\par
\noindent\hspace*{\parindent}Hätte man in obigen Beispielen $\log\operatorname{tg}i=9{,}267\,1790$. vorausgesetzt, so würde man folgende Resultate erhalten haben:
\[
\begin{array}{rl}
v-\AstNode=&12^0\;0'\,17''22\\
v=&69\;26\;47''40\\
\log\Delta=&0{,}273\,1377\\
\log r=&0{,}420\,3009
\end{array}
\]
\noindent\hspace*{\parindent}[Hier bricht die Aufzeichnung ab.]
\begin{center}[6.]\end{center}
\noindent{\small [S. 16]}\par
\noindent\hspace*{\parindent}Sämmtliche vorstehende Beispiele beziehen sich auf den neuen von \textsc{Piazzi} im Anfang dieses Jahrs zu Palermo entdeckten Planeten, und zwar auf folgende drei Beobachtungen
{\small\setlength{\arraycolsep}{3pt}
\[
\begin{array}{ll@{\quad}ll@{\quad}lll}
\text{Jan.}&2{,}360\,4699&\text{Länge}&53^0\,19'\,44''3&\text{Breite}&3^0\;2'\,24''9&\text{südlich}\\
&22{,}305\,8067&&53\;39\;1''8&&1\;42\;28''1&\\
&42{,}258\,3125&&56\;26\;40''0&&0\;36\;2''9.&
\end{array}
\]}

Die kleinen Abweichungen in obigen Annahmen [von den \textsc{Piazzi}schen Zahlen] rühren von den wegen Aberration und Präcession der Nachtgleichen
\clearpage
\HistoricalHead{240}{L}{THEORETISCHE ASTRONOMIE. NACHLASS.}
\SourceLeaf{245}
\noindent angebrachten Verbesserungen her. Zur Berechnung der Aberration sind folgende scheinbare tägliche Bewegungen gebraucht worden:
{\small\setlength{\arraycolsep}{3pt}
\[
\begin{array}{llrlr}
1^{\text{te}}\text{ Beobachtung}&\text{B[ewegung] in der Länge}&-179''85&\text{in d[er] Breite}&-256''1\\
2^{\text{te}}&&+279''2&&-216''95\\
3^{\text{te}}&&+687''5667&&-177''0667.
\end{array}
\]}
\begin{center}[7.]\end{center}
\noindent\hspace*{\parindent}Durch Interpolation [aus den mit drei verschiedenen Hypothesen über $i$ und $\AstNode$ gefundenen Resultaten] fand man sodann folgende Elemente, welche die beiden äussern Beob[achtungen] genau, die mittl[ere] mit $+1''84$ Fehler in d[er] Länge u[nd] $-2''36$ in d[er] Breite darstellen\textsuperscript{[*)]}.\begingroup\renewcommand{\thefootnote}{[*)}\footnotetext{Vgl. Werke VI, S. 200.]}\endgroup
\par\medskip\noindent{\small [S. 17]}
\[
\begin{array}{rl}
\text{Neigung der Bahn}&10^0\,32'\,19''\\
\AstNode&81\;8\;50\\
\omega&330\;14\;33\\
e=&0{,}083\,28356
\end{array}
\]
\[
\begin{array}{rl}
\text{Grösste Aequ. C[entri]}=&9^0\,32'\,57''\\
\log a=&0{,}438\,1058\\
\text{temp. rev.}=&4^a\,54\,107\\
\text{mot. diurnus}=&781''355
\end{array}
\]
\[
\text{long. med. Jan.}\;2{,}360\,4699\;\text{. . .}\;78^0\,25'\,13''.
\]
\begin{center}\rule{.25\linewidth}{.3pt}\end{center}
\clearpage
\endgroup
